\documentclass[11pt,letterpaper]{article}

% --- Paquetes de Idioma y Tipografía ---
\usepackage[utf8]{inputenc}
\usepackage[spanish,es-tabla]{babel}
\usepackage[T1]{fontenc}
\usepackage{lmodern}
\usepackage{microtype}

% --- Configuración de Página y Márgenes ---
\usepackage[top=2.3cm, bottom=2.3cm, left=2.3cm, right=2.3cm]{geometry}
\usepackage{setspace}
\setstretch{1.15}

% --- Paquetes Matemáticos y Cajas ---
\usepackage{amsmath, amssymb, amsfonts}
\usepackage[most]{tcolorbox}
\usepackage{booktabs}
\usepackage{tabularx}

% --- Enlaces e Hipervínculos ---
\usepackage{hyperref}
\hypersetup{
    colorlinks=true,
    linkcolor=blue!70!black,
    citecolor=blue!70!black,
    urlcolor=blue!70!black
}

% --- Encabezados y Pies de Página ---
\usepackage{fancyhdr}
\pagestyle{fancy}
\fancyhf{}
\fancyhead[L]{\small Computación Científica --- Primera Entrega}
\fancyhead[R]{\small Resumen y Verificación Bibliográfica}
\fancyfoot[C]{\thepage}
\renewcommand{\headrulewidth}{0.4pt}

% Definición de estilo para cajas de artículos
\newtcolorbox{paperbox}[2][]{
    enhanced,
    colback=white,
    colframe=blue!70!black,
    coltitle=white,
    fonttitle=\bfseries\small,
    title=#2,
    boxrule=0.8pt,
    arc=2mm,
    top=3mm, bottom=3mm, left=3mm, right=3mm,
    #1
}

\begin{document}

\begin{center}
    {\Large \textbf{Revisión y Resumen de la Literatura Científica Central}}\\[0.25cm]
    {\large \textbf{Sistemas Dinámicos Caóticos 1D, Coeficiente de Lyapunov y Criptografía de Imágenes}}\\[0.15cm]
    \small \textbf{Curso:} Computación Científica \quad | \quad \textbf{Pontificia Universidad Javeriana Cali}
\end{center}
\vspace{0.3cm}

\begin{tcolorbox}[colback=blue!5!white, colframe=blue!50!black, arc=1.5mm, boxrule=0.6pt]
\small
\textbf{Nota de Trazabilidad Bibliográfica:} Las cuatro fuentes centrales que sustentan teórica y metodológicamente esta Primera Entrega se derivan y conectan directamente a través del artículo general de revisión:
\begin{center}
    \textbf{Zhang, B., \& Liu, L. (2023).} \textit{Chaos-Based Image Encryption: Review, Application, and Challenges}. \\
    \textit{Mathematics} (MDPI), 11(11), 2585. DOI: \href{https://doi.org/10.3390/math11112585}{10.3390/math11112585}.
\end{center}
\end{tcolorbox}

\vspace{0.2cm}

% =========================================================================
% ARTÍCULO 1: ZHANG & LIU (2023)
% =========================================================================
\begin{paperbox}{1. Documento Matriz: Zhang \& Liu (2023)}
\textbf{Referencia completa:} Zhang, B., \& Liu, L. \textit{Chaos-based image encryption: review, application, and challenges}. \textit{Mathematics} (MDPI), 2023, 11(11), 2585. \\
\textbf{DOI:} \href{https://doi.org/10.3390/math11112585}{10.3390/math11112585} \quad | \quad \textbf{Rol:} Documento base del proyecto.
\begin{itemize}
    \item \textbf{Problema central:} La transmisión masiva y segura de imágenes digitales requiere algoritmos de cifrado de alta velocidad capaces de romper la correlación y redundancia espacial intrínseca entre píxeles adyacentes.
    \item \textbf{Aporte al proyecto:} 
    \begin{enumerate}
        \item Expone en la Sección 2.2.2 los diversos indicadores cuantitativos del caos y define en la \textbf{Ecuación (8)} el Coeficiente de Lyapunov como la medida de divergencia local: $\lambda = \lim_{n \to \infty} \frac{1}{n} \sum_{i=0}^{n-1} \ln |F'(X_i)|$, demostrando que $\lambda > 0$ es la condición requerida para seguridad criptográfica.
        \item Presenta las formulaciones del \textbf{Mapa Logístico} (Ec. 1) y del \textbf{Mapa de Tienda} (Ec. 2).
        \item Fundamenta las métricas de evaluación criptográfica (NPCR, UACI y correlación adyacente) en la Sección 4.
    \end{enumerate}
\end{itemize}
\end{paperbox}

% =========================================================================
% ARTÍCULO 2: MAY (1976)
% =========================================================================
\begin{paperbox}{2. May (1976) --- [Citado en Zhang \& Liu como Ref. 3]}
\textbf{Referencia completa:} May, R.~M. \textit{Simple mathematical models with very complicated dynamics}. \textit{Nature}, 1976, 261(5560), 459--467. \\
\textbf{DOI:} \href{https://doi.org/10.1038/261459a0}{10.1038/261459a0} \quad | \quad \textbf{Indexación:} Pág. 2 y Pág. 32 (Ref. 3).
\begin{itemize}
    \item \textbf{Problema central:} Modelado del crecimiento biológico limitado por capacidad de carga ambiental mediante la ecuación en diferencias no lineal $x_{n+1} = r x_n (1 - x_n)$.
    \item \textbf{Aporte al proyecto:} Es el artículo seminal que introdujo el \textbf{Mapa Logístico} en la literatura científica moderna. Demuestra la transición universal desde un punto fijo estable ($x^* = 1 - 1/r$ para $1 < r \le 3$), pasando por duplicaciones de período, hasta alcanzar el régimen caótico completamente desarrollado. Sustenta la clase \texttt{LogisticMap} en \texttt{src/maps.py}.
\end{itemize}
\end{paperbox}

% =========================================================================
% ARTÍCULO 3: LI & YORKE (1975)
% =========================================================================
\begin{paperbox}{3. Li \& Yorke (1975) --- [Citado en Zhang \& Liu como Ref. 2]}
\textbf{Referencia completa:} Li, T.-Y., \& Yorke, J.~A. \textit{Period three implies chaos}. \textit{The American Mathematical Monthly}, 1975, 82(10), 985--992. \\
\textbf{DOI:} \href{https://doi.org/10.2307/2318254}{10.2307/2318254} \quad | \quad \textbf{Indexación:} Pág. 2 y Pág. 32 (Ref. 2).
\begin{itemize}
    \item \textbf{Problema central:} Determinar bajo qué condiciones topológicas una transformación continua unidimensional exhibe órbitas aperiódicas complejas.
    \item \textbf{Teorema fundamental:} Demuestra formalmente que la existencia de una órbita periódica de período 3 implica la existencia de órbitas periódicas de cualquier longitud y un conjunto no numerable de condiciones iniciales caóticas.
    \item \textbf{Aporte al proyecto:} Justifica matemáticamente la aparición de la \textbf{ventana periódica de período 3} en el Mapa Logístico ($r \approx 3.8284$), explicando por qué el Coeficiente de Lyapunov cae súbitamente a valores negativos ($\lambda < 0$) en medio de la región caótica.
\end{itemize}
\end{paperbox}

% =========================================================================
% ARTÍCULO 4: AMIN ET AL. (2010)
% =========================================================================
\begin{paperbox}{4. Amin, Faragallah, \& Abd El-Latif (2010) --- [Citado en Zhang \& Liu como Ref. 16]}
\textbf{Referencia completa:} Amin, M., Faragallah, O.~S., \& Abd El-Latif, A.~A. \textit{A chaotic block cipher algorithm for image cryptosystems}. \textit{Communications in Nonlinear Science and Numerical Simulation}, 2010, 15(11), 3484--3497. \\
\textbf{DOI:} \href{https://doi.org/10.1016/j.cnsns.2009.12.016}{10.1016/j.cnsns.2009.12.016} \quad | \quad \textbf{Indexación:} Pág. 4 (Ec. 2) y Pág. 33 (Ref. 16).
\begin{itemize}
    \item \textbf{Problema central:} Diseñar un algoritmo de cifrado para imágenes rápido y robusto basado en mapas unidimensionales, superando las debilidades del mapa logístico (ventanas periódicas) mediante el Mapa de Tienda y difusión encadenada.
    \item \textbf{Aporte al proyecto:} 
    \begin{enumerate}
        \item Define el \textbf{Mapa de Tienda (\textit{Tent Map})} y demuestra que su Coeficiente de Lyapunov es estrictamente monótono $\lambda(p) = \ln(p)$, garantizando caos continuo sin colapso periódico para todo $p > 1$.
        \item Fundamenta el criptosistema implementado en \texttt{src/cipher.py}: confusión mediante ordenamiento caótico y difusión XOR encadenada $C_i = P'_i \oplus S_i \oplus C_{i-1}$, permitiendo descifrado exacto y resistencia a ataques diferenciales.
    \end{enumerate}
\end{itemize}
\end{paperbox}

\end{document}
